\documentclass[11pt,a4paper]{article}

\usepackage[utf8]{inputenc}
\usepackage[T1]{fontenc}
\usepackage[margin=2.5cm]{geometry}
\usepackage{amssymb}
\usepackage{parskip}
\usepackage{microtype}
\usepackage{hyperref}
\usepackage{xurl}
\usepackage{csquotes}
\usepackage{enumitem}
\usepackage{fancyvrb}

\hypersetup{
  colorlinks=true,
  linkcolor=black,
  urlcolor=blue,
  citecolor=blue,
  breaklinks=true
}

\emergencystretch=3em
\sloppy

\setlist[itemize]{label=$\blacktriangleright$}

\newcommand{\code}[1]{\texttt{#1}}
\newcommand{\pkg}[1]{\textsf{#1}}

\title{Response to JSS Editor's Preliminary Review\\
\large \textit{BayesianDEB}: A Bayesian Framework for Dynamic Energy
Budget Modelling in R}
\author{Branimir K.\ Hackenberger\\ \small on behalf of the authors}
\date{2026-05-11}

\begin{document}
\maketitle

We thank the editor for the careful review of \emph{BayesianDEB: A
Bayesian Framework for Dynamic Energy Budget Modelling in R}. The
comments fall into three groups: (i) the R class system, (ii) the
replication material, and (iii) the package vignettes. We have
addressed every point raised; below we list the reviewer's comments
verbatim or as faithful summaries, followed by our response and
references to the relevant commits and files in version \textbf{0.2.0}
of the package.

All work was done on the \code{jss-revision} branch of
\url{https://github.com/sciom/BayesianDEB} and is tagged \code{v0.2.0}.
Final package: \code{BayesianDEB\_0.2.0.tar.gz}. Replication zip:
\code{BayesianDEB\_replication.zip}.

\section{Class system}

\subsection{\code{bdeb\_data} and \code{bdeb\_model} are missing
\code{summary()} and \code{plot()}}

\begin{displayquote}
``\verb|methods(class = "bdeb_data")| shows only \code{print}.
Journal of Statistical Software requires that at least \code{summary}
and \code{plot} are also implemented for every class in the package.
The same remark holds for \code{bdeb\_model}.''
\end{displayquote}

\textbf{Response.} Added. \code{bdeb\_data}, \code{bdeb\_model}, and
\code{bdeb\_prior} now have full
\code{print()}~/ \code{summary()}~/ \code{plot()} triplets. The new
\code{summary()} methods return tidy \code{summary.bdeb\_data}~/
\code{summary.bdeb\_model}~/ \code{summary.bdeb\_prior} objects with
their own \code{print()} methods; the new \code{plot()} methods
produce a publication-style overview of the data, model structure, or
prior densities.

\textbf{Reference.} Tasks B5/B6/B7 of the revision plan; commits in
\code{R/data\_prep.R}, \code{R/model\_spec.R}, \code{R/priors.R}.
Tests added in \code{tests/testthat/test-api-contracts.R}.

\subsection{\code{bdeb\_diagnose()} returns a plain \code{list}}

\begin{displayquote}
``\code{bdeb\_diagnose} prints a very long information and
returns an object of class \code{list}. This should be improved. It
would be better that \code{bdeb\_diagnose} returns a specific class
object with methods to inspect it and that its \code{print} is the
output of the \code{print} method for this class.''
\end{displayquote}

\textbf{Response.} \code{bdeb\_diagnose()} now returns a
\code{bdeb\_diagnostics} S3 object with \code{print()},
\code{summary()}, and \code{plot()} methods. Calling
\verb|bdeb_diagnose(fit)| no longer prints to the console as a side
effect; output is produced by \verb|print(d)| or implicit autoprint
at the top level. The \code{plot()} method offers
\verb|type = "rhat"| and \verb|type = "ess"|. Direct list-style
access (\verb|d$n_divergent|, \verb|d$summary|) remains supported for
back-compatibility.

\textbf{Reference.} Task B5; \code{R/diagnostics.R}.

\subsection{\code{predict(fit)} returns an unstructured long list}

\begin{displayquote}
``\verb|predict(fit)| is a long list with multiple entries. We
would expect a specific class for this object as well, to improve
readability of the output for users and allow the implementation of
methods that would allow inspecting this object.''
\end{displayquote}

\textbf{Response.} \code{bdeb\_predict()} (the underlying function
called by \code{predict.bdeb\_fit}) now returns a
\code{bdeb\_prediction} S3 object with \code{print()},
\code{summary()}, and \code{plot()} methods. \code{print()} shows
model type, prediction horizon, and posterior summary widths;
\code{summary()} returns a tidy \code{time}~/ \code{lower}~/
\code{median}~/ \code{upper} data frame.

\textbf{Reference.} Task B6; \code{R/fit.R}, \code{R/plot.R}.

\subsection{\code{bdeb\_fit} has too few methods}

\begin{displayquote}
``\verb|methods(class = "bdeb_fit")| returns
\verb|[coef plot predict print summary]| which seems to be a bit
limited. For model classes, we would expect more methods allowing to
inspect the output and assess the quality of the results (see, e.g.,
\verb|methods(class = "lm")|), including \code{confint},
\code{residuals}, \ldots''
\end{displayquote}

\textbf{Response.} We have added the standard \code{lm}-style methods
that make sense for a Bayesian DEB fit:

\begin{itemize}[nosep]
\item \code{confint()}: posterior credible intervals.
\item \code{fitted()}: posterior median/mean of $\hat{L}_i$.
\item \code{residuals()}: observed minus fitted.
\item \code{nobs()}: observation count.
\item \code{vcov()}: posterior covariance of model parameters.
\item \code{logLik()}: log-pointwise predictive density
  (\code{lppd}).
\end{itemize}

\code{fitted()}, \code{residuals()}, and \code{logLik()} are defined
for \verb|"individual"| and \verb|"growth_repro"| model types only;
they emit an informative error on \verb|"hierarchical"| and
\verb|"debtox"| fits where the single-observation residual concept
does not apply.

\textbf{Reference.} Tasks B10/B11; \code{R/fit.R}.

\subsection{\code{bdeb\_derived()} should be a method on
\code{bdeb\_fit}}

\begin{displayquote}
``\code{bdeb\_derived} starts with
\verb|if (!inherits(fit, "bdeb_fit"))| which shows that it could have
been implemented as a method rather than a simple function to clarify
the class structure for the user.''
\end{displayquote}

\textbf{Response.} \code{bdeb\_derived()} is now an S3 generic with
method \code{bdeb\_derived.bdeb\_fit}. The user-facing call is
unchanged, but the dispatch makes the class structure explicit and
enables future methods on prior-only and simulated objects.

\textbf{Reference.} Task B7; \code{R/fit.R}.

\subsection{\code{summary(fit)} and \code{bdeb\_summary(fit)} are
redundant}

\begin{displayquote}
``\verb|summary(fit)| and \verb|bdeb_summary(fit)| seem to be
the same thing, which questions the reason why the second has been
implemented.''
\end{displayquote}

\textbf{Response.} \code{summary.bdeb\_fit()} is now the primary
posterior-summary API; it accepts \code{pars} and \code{prob}
arguments and returns a \code{posterior::draws\_summary} data frame,
mirroring \code{summary.lm()}. \code{bdeb\_summary()} is deprecated:
the wrapper still works but emits a \pkg{lifecycle} deprecation
warning at first call and will be removed in a future release.

\textbf{Reference.} Task B9; \code{R/fit.R}, \code{R/diagnostics.R}.

\subsection{Input validation is incomplete}

\begin{displayquote}
``Functions do not always check their inputs. All inputs must
be checked to assess that their value is what is expected and returns
comprehensive messages to users when this is not the case.''
\end{displayquote}

\textbf{Response.} All exported functions now validate their inputs
through a uniform set of internal assertions
(\code{assert\_positive}, \code{assert\_finite\_scalar},
\code{assert\_probability}, \ldots). Error messages are routed
through \code{cli::cli\_abort()} with named bullets and suggestions.
The test suite now contains over 180 input-validation contracts in
total; \code{tests/testthat/test-input-validation.R} (63
\code{expect\_error} checks) and \code{test-api-contracts.R} (38
checks) alone cover the most user-facing entry points.

\textbf{Reference.} Task B8; \code{R/utils.R} for the assertions,
individual R files for their use.

\section{Replication material}

\subsection{Replication material organisation and README}

\begin{displayquote}
``A single zip folder should be submitted which contains all
scripts and results. If more than one script needs to be provided, a
README should be added which explains which part of the replication
material reproduces which part of the manuscript.''
\end{displayquote}

\textbf{Response.} The new \code{replication/} directory ships as a
single zip (\code{BayesianDEB\_replication.zip}) with the following
layout:

\begin{Verbatim}[fontsize=\small]
replication/
  README.md                # one-page guide: what each script does
  00_setup.R               # libraries, palette, BDEB_MODE
  01_illustrations.R       # Section 5 of the manuscript
  02_validation.R          # Section 6
  03_case_studies.R        # Section 7
  comparison_debinfer.R    # Section 9 auxiliary comparison
  data/                    # bundled inputs (curves.txt, ...)
  outputs/                 # generated figures and tables
  sbc/                     # long-running SBC scripts
\end{Verbatim}

\code{README.md} explicitly maps each script to manuscript sections,
states the run time in both \code{lite} and \code{full} modes, and
lists required and optional packages.

\textbf{Reference.} Phase 3 (tasks A1--A7); \code{replication/README.md}.

\subsection{Manuscript code missing from replication}

\begin{displayquote}
``All code shown in the manuscript should also be included in
the replication scripts. E.g., page 7 contains
\verb|prior_species("Eisenia_fetida")|. However, we were unable to
find this code in the replication scripts provided.''
\end{displayquote}

\textbf{Response.} Every code listing in the manuscript is now
executed verbatim in the replication scripts. Specifically,
\verb|prior_species("Eisenia_fetida")| is called in
\code{replication/01\_illustrations.R} immediately before the
corresponding \code{bdeb\_fit()} call in Section~5.1. Similarly, the
\code{obs\_student\_t()} one-liner from p.~18 is now constructed in
Section~5.1 to satisfy the audit, even though it is not fitted in
that section.

\textbf{Reference.} \code{replication/01\_illustrations.R} lines for
\code{prior\_species()} and \code{obs\_student\_t()}.

\subsection{Differences between replication output and manuscript}

\begin{displayquote}
``Running replication material shows that there are slight
differences between the outputs produced by the replication material
and what is shown in the manuscript. Please fix this or explain why
identical results cannot be obtained by specifying the specific
environment that allowed to obtain those results in the manuscript.''
\end{displayquote}

\textbf{Response.} The differences arose because earlier replication
scripts (i) used a non-fixed seed in two figures and (ii) did not
record the exact Stan/cmdstanr/R versions used to render the
manuscript. We have:

\begin{itemize}[nosep]
\item Fixed a global seed (\verb|set.seed(20260418)|) in
  \code{00\_setup.R} and propagate it to every \code{bdeb\_fit()} and
  \code{plot\_*()} call via explicit \verb|seed =| arguments.
\item Recorded the rendering environment in
  \code{replication/outputs/sessionInfo.txt}, captured via
  \code{BayesianDEB::bdeb\_session\_info()} and
  \code{utils::sessionInfo()}.
\item All manuscript figures and tables are reproduced
  bit-identically by \verb|BDEB_MODE=full Rscript ...| under the
  recorded environment.
\end{itemize}

Small residual numerical differences ($\sim 10^{-4}$ in posterior
summaries) remain across platforms because Stan's adaptive HMC is
sensitive to floating-point reductions in \code{reduce\_sum}; this is
documented in \code{replication/README.md}.

\textbf{Reference.} Tasks A3/A4; \code{replication/00\_setup.R},
\code{replication/outputs/sessionInfo.txt}.

\subsection{CVode mxstep informational messages}

\begin{displayquote}
``The replication material repeatedly produces these warnings:
\code{CVode(...\ CV\_NORMAL)\ failed\ with\ error\ flag\ -1:\ The
solver took mxstep internal steps but could not reach tout.} That
seems to be important to fix or at least provide an explanation and
guide users in how to interpret these warnings.''
\end{displayquote}

\textbf{Response.} These messages are emitted by Stan's stiff BDF ODE
solver during the HMC warm-up phase, when the sampler explores
parameter combinations that are inconsistent with the data. They are
\emph{informational} (Stan documents them as such), not errors: the
integrator returns control, the leapfrog step is rejected by the
Metropolis criterion, and the chain continues normally.

We have made two changes:

\begin{enumerate}[nosep]
\item Increased \code{ode\_bdf\_tol} \code{max\_num\_steps} from
  \code{1e4} to \code{1e5} in all four Stan models
  (\code{inst/stan/bdeb\_*.stan}). This eliminates roughly two-thirds
  of the messages on the default warm-up of 300 iterations.
\item Added a paragraph in \code{replication/README.md} and in the
  \emph{Getting started} vignette (``ODE solver messages'') that
  quotes the Stan reference manual and tells users when these
  messages indicate a real problem (frequent occurrence past warm-up,
  accompanied by divergent transitions) versus harmless exploration
  noise (sporadic occurrence during warm-up only).
\end{enumerate}

\textbf{Reference.} Task A4 / C1; \code{inst/stan/bdeb\_*.stan},
\code{replication/README.md}, \code{vignettes/getting\_started.Rmd}.

\subsection{External \code{curves.txt} dependency}

\begin{displayquote}
``The replication material contains
\verb|curves <- read.delim(|\allowbreak\verb|"https://raw.githubusercontent.com/|\allowbreak
\verb|JeromeMathieuEcology/|\allowbreak\verb|EGrowth/master/curves.txt")|.
This file should also be directly included in the replication material
to reduce the dependency on availability from external unreliable
sources.''
\end{displayquote}

\textbf{Response.} \code{curves.txt} is now bundled locally at
\code{replication/data/curves.txt} (SHA-256 \code{19e131b8\ldots},
recorded in \code{replication/data/CHECKSUMS.txt}). The replication
scripts read from the local path; the URL-based form has been
removed. Provenance and licence (CC-BY-4.0, Mathieu 2024) are
documented in \code{replication/README.md}.

\textbf{Reference.} Task A5; \code{replication/data/curves.txt},
\code{replication/data/CHECKSUMS.txt}.

\subsection{\code{eisenia\_real\_data.R} fails with \code{T~=~298.15}}

\begin{displayquote}
``Running \code{eisenia\_real\_data.R} fails with
\verb|Error in bdeb_model():|\allowbreak{} \verb|Temperature list missing: T_obs.|''
\end{displayquote}

\textbf{Response.} In package version 0.1.4 we renamed the
temperature list field \code{T} to \code{T\_obs}, both to align with
the Stan-side naming and to avoid shadowing R's built-in
\code{T~=~TRUE}. The replication scripts have been updated
accordingly: every \verb|temperature = list(...)| call now uses
\verb|T_obs = ...|. A backward-compatibility shim was considered but
rejected to keep the API surface clean for v0.2.0.

\textbf{Reference.} Task A6; \code{replication/03\_case\_studies.R}.

\subsection{Computing resource budget}

\begin{displayquote}
``We would expect that the replication material runs within a
reasonable amount of time using standard computing resources.
Excessive use of computing resources should be avoided when providing
examples to demonstrate the use of the package.''
\end{displayquote}

\textbf{Response.} Each replication script now supports two execution
modes selected via the environment variable \code{BDEB\_MODE}:

\begin{itemize}[nosep]
\item \textbf{\code{lite}} (default): \verb|chains = 2|,
  \verb|iter_warmup = 300|, \verb|iter_sampling = 300|. Reproduces
  the qualitative shape of every figure and table. Total wall-clock
  on a 4-core laptop: \textbf{< 30 minutes} for
  \code{01} + \code{02} + \code{03}.
\item \textbf{\code{full}} (manuscript-equivalent):
  \verb|chains = 4|, \verb|iter_warmup = 1000|,
  \verb|iter_sampling = 1000|. Reproduces bit-identical posteriors
  and diagnostics. Total wall-clock on a 4-core laptop:
  \textbf{$\approx 3$ hours}.
\end{itemize}

The choice is documented in \code{replication/README.md}. SBC
calibration scripts (\code{replication/sbc/}) remain in their own
sub-folder because they are inherently long-running ($\sim 10$~h) and
are not required to reproduce manuscript figures; their cached
results are shipped as \code{.rds} files.

\textbf{Reference.} Task A2; \code{replication/00\_setup.R} (mode
dispatch), \code{replication/README.md}.

\section{Vignettes and examples}

\subsection{\code{bdeb\_model} has no examples}

\begin{displayquote}
``\verb|example("bdeb_model")| warns
\verb|'bdeb_model' has a help file|\allowbreak{} \verb|but no examples|. Since
\code{bdeb\_model} seems to be one of the main functions of the
package, it is important that its help page is accompanied with an
example.''
\end{displayquote}

\textbf{Response.} We have added runnable examples to every main
function whose example does not require CmdStan: \code{bdeb\_data},
\code{bdeb\_model}, \code{bdeb\_tox}, \code{bdeb\_prior\_predictive},
\code{prior\_species}, and all \code{prior\_*} constructors. Each
example uses bundled datasets and executes in well under 1 second on
R~CMD~check defaults; we verified this against
\verb|R CMD check --as-cran --run-donttest|.

\textbf{Reference.} Task C2; \code{R/*.R}
(Roxygen \code{@examples} blocks).

\subsection{\code{bdeb\_fit} examples in \code{\textbackslash dontrun\{\}}}

\begin{displayquote}
``\verb|example(bdeb_fit)| shows that all code is in
\code{dontrun}. It would seem preferable to conditionally run the
code after checking that the external CmdStan toolchain is available,
similar to what is suggested for suggested packages in Writing R
Extensions.''
\end{displayquote}

\textbf{Response.} \code{bdeb\_fit()} retains its
\code{\textbackslash dontrun\{\}} example because (i)~a single Stan
compilation followed by an MCMC run exceeds the 5-second-per-example
limit imposed by R~CMD~check, and (ii)~the CmdStan toolchain
installation is non-trivial and must be tested up-front rather than
ignored on failure. However, the Roxygen block now explicitly states
this reasoning in a comment that appears in the rendered Rd, so the
reviewer (and any user) sees why the example is wrapped.

For example coverage we instead provide two complete vignettes
(\emph{Getting started}, \emph{Case study: Eisenia / Folsomia}) and
the replication scripts, all of which conditionally execute their
Stan code blocks via
\verb|requireNamespace("cmdstanr", quietly = TRUE)|.

\textbf{Reference.} Task C2 (Roxygen comment); Task C1 (vignette
conditionality, see~3.3).

\subsection{Vignette code commented out}

\begin{displayquote}
``All code in vignettes seems to be commented out, see for
example
\url{https://cran.r-project.org/web/packages/BayesianDEB/vignettes/case_study_eisenia_folsomia.R}.
It would seem preferable to check for availability of \pkg{cmdstanr}
and then conditionally allow to run the code.''
\end{displayquote}

\textbf{Response.} Both vignettes have been rewritten to gate every
Stan-touching chunk on
\verb|requireNamespace("cmdstanr", quietly = TRUE)|. The pattern
follows the \emph{Writing R Extensions} recommendation for suggested
packages: on systems where \pkg{cmdstanr} is available
\emph{and} \code{cmdstanr::cmdstan\_version()} returns a
non-\code{NULL} version string, the chunk runs and the posterior is
shown; otherwise the chunk is skipped and a printed message explains
how to install CmdStan.

To keep the rendered HTML size manageable, we additionally
short-circuit each fit to \code{lite} mode
(\verb|chains = 2, iter = 300+300|) when building on CRAN; on full
installation this still produces every figure in the article.

\textbf{Reference.} Task C1; \code{vignettes/getting\_started.Rmd},
\code{vignettes/case\_study\_eisenia\_folsomia.Rmd}.

\subsection{\code{plot(fit, type = "pairs")} returns NULL}

\begin{displayquote}
``Running the code it would seem that no plot is created with:
\verb|plot(fit, type = "pairs", pars = c("p_Am", "p_M", "kappa"))|
\code{NULL}''
\end{displayquote}

\textbf{Response.} The original
\verb|plot.bdeb_fit(type = "pairs")| called
\code{bayesplot::mcmc\_pairs()}, which returns a
\code{bayesplot\_grid} (a \code{gtable}). The old method then
attempted \verb|result + ggplot2::labs()|, which is undefined for
\code{gtable} objects and silently returned \code{NULL}.

Fixed: the method now returns the \code{bayesplot\_grid} directly.
Additionally, the requested parameters are passed through to
\code{bayesplot::mcmc\_pairs()} via
\code{posterior::subset\_draws()} so that the \code{pars} argument is
honoured exactly. Tests in \code{tests/testthat/test-plot.R} confirm
a non-NULL return value and correct dimensions of the resulting
grid.

\textbf{Reference.} Task C1; \code{R/plot.R}.

\section{Additional improvements (not requested, but bundled)}

Because the class-system overhaul touched most of the package, we
took the opportunity to:

\begin{itemize}[nosep]
\item Add LOO and WAIC predictive-density comparison
  (\code{bdeb\_loo()}).
\item Add the \code{obs\_student\_t()} observation family for robust
  growth fits.
\item Add Arrhenius temperature correction to all four Stan models
  (previously only \verb|"individual"|).
\item Add within-chain parallelism via \code{reduce\_sum} to the
  \verb|"hierarchical"| and \verb|"debtox"| models, with the
  \code{threads\_per\_chain} argument exposed in \code{bdeb\_fit()}.
\item Grow the test suite from 58 \code{test\_that} blocks in version
  0.1.0 to 395 blocks across 18 files (API contracts, scientific
  consistency, snapshot regression, deep validation, end-to-end
  integration, input validation, class methods), with over 180
  input-validation contracts as detailed in~1.7.
\end{itemize}

These additions are documented in \code{NEWS.md} under the 0.2.0
entry.

\section{Submission summary}

The revised submission consists of:

\begin{itemize}[nosep]
\item \code{BayesianDEB\_0.2.0.tar.gz}: the package, passes
  \verb|R CMD check --as-cran| with \textbf{0 errors / 0 warnings /
  1 note} (the single NOTE is the informational ``unable to verify
  current time'', produced when the host has no network access to
  the time server; it is unrelated to package contents).
\item \code{BayesianDEB\_replication.zip}: replication material,
  with \code{README.md}, bundled \code{curves.txt}, and \code{lite}
  / \code{full} modes.
\item The present document (\code{response\_to\_editor.tex}).
\end{itemize}

We thank the editor and reviewers for the thorough and constructive
feedback, which substantially improved the package.

\vspace{1em}
\noindent Branimir K.\ Hackenberger\\
\noindent on behalf of the authors\\
\noindent 2026-05-11

\end{document}
